\section{Fundamentação Teórica e Trabalhos Relacionados}

\subsection{Cinemática Diferencial e Odometria (RK2)}

A fim de modelar a cinemática de um robô móvel, o processo ocorre de
baixo para cima (\textit{bottom-up}), pois cada roda contribui         
individualmente para o movimento e também impõe restrições ao    
deslocamento do chassi~\cite{siegwart2004}. Desse modo, as forças e      
restrições de cada roda devem ser expressas em relação a um referencial  
consistente, tornando essencial o mapeamento entre o referencial global  
$\{X_I, Y_I\}$ e o referencial local $\{X_R, Y_R\}$, fixado no ponto     
médio $P$ do eixo motriz~\cite{siegwart2004}. Considerando a pose planar
$\boldsymbol{\xi}_I = [x,\, y,\, \theta]^T$, o mapeamento do movimento do
referencial global para o referencial local é dado
por~\cite{siegwart2004}:
\begin{equation}
\dot{\boldsymbol{\xi}}_R = R(\theta)\,\dot{\boldsymbol{\xi}}_I, 
\qquad
R(\theta) =
\begin{bmatrix}
\cos\theta & \sin\theta & 0 \\
-\sin\theta & \cos\theta & 0 \\ 
0 & 0 & 1 
\end{bmatrix}, 
\label{eq:rotacao_global_local}
\end{equation} 
em que $R(\theta)$ é a matriz de rotação ortogonal, $\theta$ representa a
orientação do robô no plano e $\dot{\boldsymbol{\xi}}$ representa o vetor
de velocidades (derivada da pose). Para mapear as velocidades do 
referencial local para o referencial global, utiliza-se a matriz inversa 
de $R(\theta)$: $\dot{\boldsymbol{\xi}}_I =
R(\theta)^{-1}\dot{\boldsymbol{\xi}}_R =
R(\theta)^T\dot{\boldsymbol{\xi}}_R$~\cite{siegwart2004}. Como nenhuma 
das rodas permite deslocamento lateral sem escorregamento, a velocidade 
lateral local é nula ($\dot{y}_R = 0$), restando apenas a velocidade 
longitudinal $\dot{x}_R = v$ e a velocidade angular $\dot{\theta}_R =
\omega$. Considerando o raio das rodas $R$ e a bitola $L = 2l$ entre     
elas, a cinemática direta a partir das velocidades angulares das rodas
direita ($\omega_d$) e esquerda ($\omega_e$) é dada 
por~\cite{siegwart2004}: 
\begin{equation}
v = \frac{R}{2}(\omega_d + \omega_e), \qquad 
\omega = \frac{R}{L}(\omega_d - \omega_e),
\label{eq:cinematica_direta} 
\end{equation}
resultando no referencial global em $\dot{x} = v\cos\theta$, $\dot{y} =
v\sin\theta$ e $\dot{\theta} = \omega$.

Conforme o robô se movimenta no ambiente, erros sistemáticos e aleatórios
nos \textit{encoders} e no contato com o solo fazem a odometria acumular 
erro de posição e, principalmente, deriva angular entre o referencial 
estimado e o original~\cite{siegwart2004,corke2023}. Na integração 
numérica discreta ao longo de um intervalo $\Delta t$, sendo $\Delta s_k 
= v_k\,\Delta t$ e $\Delta\theta_k = \omega_k\,\Delta t$, o método
clássico de Runge-Kutta de 2ª~ordem (RK2, ou Método do Ponto Médio) reduz
o erro de truncamento em curvas em comparação ao método de Euler        
explícito. Enquanto o método de Euler assume que o robô percorre toda a
distância $\Delta s_k$ mantendo a orientação inicial $\theta_k$ e gira
apenas ao final do passo~\cite{corke2023}, o RK2 considera que a       
orientação varia gradualmente de $\theta_k$ até $\theta_k +     
\Delta\theta_k$ durante $\Delta t$, adotando a direção média $\theta_k +
\Delta\theta_k/2$~\cite{siegwart2004}:
\begin{align}
x_{k+1} &= x_k + \Delta s_k \cos\!\left(\theta_k +
\frac{\Delta\theta_k}{2}\right), \label{eq:rk2_x} \\
y_{k+1} &= y_k + \Delta s_k \sin\!\left(\theta_k +
\frac{\Delta\theta_k}{2}\right), \label{eq:rk2_y} \\
\theta_{k+1} &= \operatorname{wrapToPi}(\theta_k + \Delta\theta_k),     
\label{eq:rk2_theta}
\end{align}
em que $\operatorname{wrapToPi}(\alpha) =
\operatorname{atan2}(\sin\alpha,\, \cos\alpha)$ mantém a orientação 
normalizada no intervalo $(-\pi, \pi]$.

\subsection{Campos Potenciais Artificiais (APF) e Limitações}

Para conduzir o robô até uma meta desviando de obstáculos, uma abordagem
clássica é o método dos Campos Potenciais Artificiais
(APF)~\cite{khatib1986,siegwart2004}. Nessa formulação, o robô move-se 
sob a influência de um campo potencial artificial $U(\mathbf{q})$, no 
qual a meta exerce atração e os obstáculos exercem
repulsão~\cite{siegwart2004}: 
\begin{align} 
U(\mathbf{q}) &= U_{\text{att}}(\mathbf{q}) + U_{\text{rep}}(\mathbf{q}),
\label{eq:apf_u_total} \\
U_{\text{att}}(\mathbf{q}) &= \frac{1}{2}\, k_{\text{att}}\,
\rho_{\text{goal}}^2(\mathbf{q}), \label{eq:apf_u_att} \\
U_{\text{rep}}(\mathbf{q}) &=
\begin{cases} 
\dfrac{1}{2}\, k_{\text{rep}} \left( \dfrac{1}{\rho(\mathbf{q})} -     
\dfrac{1}{\rho_0} \right)^2, & \text{se } \rho(\mathbf{q}) \le \rho_0, \\
0, & \text{se } \rho(\mathbf{q}) > \rho_0,
\end{cases} \label{eq:apf_u_rep} 
\end{align} 
em que $\mathbf{q} = [x, y]^T$ é a posição do robô, $U(\mathbf{q})$ é o
campo potencial resultante, $U_{\text{att}}(\mathbf{q})$ é o potencial
atrativo definido em~\eqref{eq:apf_u_att}, sendo $k_{\text{att}} > 0$
o ganho de atração e $\rho_{\text{goal}}(\mathbf{q}) = \|\mathbf{q} -  
\mathbf{q}_{\text{goal}}\|$ a distância euclidiana até a meta e    
$U_{\text{rep}}(\mathbf{q})$ é o potencial repulsivo definido            
em~\eqref{eq:apf_u_rep}, em que $k_{\text{rep}} > 0$ é o ganho de     
repulsão, $\rho(\mathbf{q})$ é a menor distância ao obstáculo e $\rho_0$
é o raio de influência repulsiva~\cite{siegwart2004}. 

Em termos de forças, o gradiente negativo do potencial fornece      
$\mathbf{F}(\mathbf{q}) = -\nabla U(\mathbf{q}) =
\mathbf{F}_{\text{att}}(\mathbf{q}) +
\mathbf{F}_{\text{rep}}(\mathbf{q})$~\cite{siegwart2004}:
\begin{align}
\mathbf{F}_{\text{att}}(\mathbf{q}) &= -k_{\text{att}}\,(\mathbf{q} -    
\mathbf{q}_{\text{goal}}), \label{eq:apf_f_att} \\
\mathbf{F}_{\text{rep}}(\mathbf{q}) &= 
\begin{cases} 
k_{\text{rep}} \left( \dfrac{1}{\rho(\mathbf{q})} - \dfrac{1}{\rho_0} 
\right) \dfrac{\mathbf{q} - \mathbf{q}_{\text{obs}}}{\rho^3(\mathbf{q})}, & \text{se } \rho(\mathbf{q})
\le \rho_0, \\
\mathbf{0}, & \text{se } \rho(\mathbf{q}) > \rho_0,
\end{cases} \label{eq:apf_f_rep}
\end{align}
em que $\mathbf{q}_{\text{obs}}$ denota o ponto mais próximo do obstáculo, considerando-se a superposição vetorial das forças repulsivas quando   
múltiplos obstáculos ou pontos detectados estão dentro do raio
$\rho_0$~\cite{khatib1986}. A magnitude e a direção de          
$\mathbf{F}(\mathbf{q})$ definem o movimento desejado do robô em direção 
à meta~\cite{siegwart2004}.

Contudo, o APF apresenta limitações inerentes, como a formação de mínimos
locais em obstáculos côncavos em formato de ``U'', nos quais as forças   
atrativa e repulsivas se anulam mutuamente, resultando em     
$\mathbf{F}(\mathbf{q}) = \mathbf{0}$ e impedindo o robô de sair da 
armadilha~\cite{khatib1986,koren1991}. Além dessa limitação, destacam-se
a dificuldade de atravessar passagens estreitas entre obstáculos 
próximos, o surgimento de oscilações em corredores e na proximidade de 
obstáculos~\cite{koren1991}, e o fenômeno GNRON (\textit{Goals
Nonreachable with Obstacles Nearby}), que ocorre quando um obstáculo está
muito próximo da meta: à medida que o robô se aproxima do alvo, a força
atrativa diminui enquanto a força repulsiva aumenta, fazendo com que a
posição da meta deixe de ser o mínimo global do potencial total e
impedindo o robô de alcançá-la~\cite{ge2000}.


\subsection{Algoritmo Bug2}

O algoritmo Bug2 é um método de navegação reativa que conduz o robô do
ponto de partida $S$ até o alvo $T$ combinando deslocamento em linha reta
com contorno de obstáculos~\cite{lumelsky1987}. Inicialmente, traça-se   
uma reta imaginária ligando $S$ a $T$, denominada linha-$m$              
(\textit{m-line})~\cite{corke2023}. O robô segue pela linha-$m$ em
direção ao alvo até encontrá-lo ou colidir virtualmente com um obstáculo,
registrando esse local como ponto de impacto $H^j$~\cite{lumelsky1987}.  
Ao detectar o obstáculo, o robô passa a contornar sua borda em uma       
direção fixa até reencontrar a linha-$m$ em um ponto $Q$ com caminho     
livre para o alvo que satisfaça a condição de progresso:                 
\begin{equation}
\|Q - T\| < \|H^j - T\|, 
\label{eq:bug2_progresso}
\end{equation}
em que $\|Q - T\|$ e $\|H^j - T\|$ representam as distâncias euclidianas
dos pontos $Q$ e $H^j$ até o alvo $T$. Quando essa condição é atendida, o
ponto $Q$ é definido como ponto de saída $L^j$ e o robô retoma o percurso
pela linha-$m$ rumo a $T$; caso o contorno complete uma volta fechada
retornando a $H^j$ sem encontrar saída válida, conclui-se que o alvo é   
inatingível~\cite{lumelsky1987}.

Sob condições ideais, o Bug2 garante convergência em tempo finito sempre
que houver caminho livre até a meta~\cite{lumelsky1987}. Contudo, em   
robôs móveis diferenciais reais ou simulados, a estratégia não considera 
restrições não-holonômicas nas mudanças de direção e apresenta
sensibilidade a ruídos dos sensores de alcance e à deriva odométrica ao 
tentar retornar com precisão à linha-$m$~\cite{siegwart2004,corke2023}.


\subsection{Trabalhos Relacionados}

Diversos estudos clássicos investigaram estratégias de navegação reativa sem mapa global prévio. Khatib~\cite{khatib1986} formulou o método dos Campos Potenciais Artificiais para desvio de obstáculos em tempo real combinando atração à meta e repulsão local, enquanto Koren e Borenstein~\cite{koren1991} analisaram suas limitações inerentes, destacando a ocorrência de mínimos locais em obstáculos côncavos, bloqueios em passagens estreitas e oscilações, ao passo que Ge e Cui~\cite{ge2000} formalizaram o problema GNRON quando há obstáculos próximos à meta. Como alternativa baseada em contorno topológico, Lumelsky e Stepanov~\cite{lumelsky1987} propuseram o algoritmo Bug2, demonstrando convergência em tempo finito mesmo diante de obstáculos desconhecidos de forma arbitrária. De modo complementar, Siegwart e Nourbakhsh~\cite{siegwart2004} e Corke~\cite{corke2023} sistematizam a modelagem cinemática diferencial, o acúmulo de erros na odometria discreta e os aspectos práticos desses métodos reativos em robôs móveis. Diferentemente de análises puramente teóricas ou com pose ideal, o presente trabalho avalia comparativamente o APF e o Bug2 sob odometria RK2 com ruído controlado, investigando o compromisso entre suavidade de trajetória e convergência em cenários com obstáculos esparsos, armadilha em ``U'' e bloqueio próximo à meta.
